\section{The End of All our Exploring}

Musings on excellence, the purpose of life, and the dissident right.

\begin{quotex}
\begin{verse}
We shall not cease from exploration\\ And the end of all our exploring\\ Will be to arrive where we started\\ And know the place for the first time.
\end{verse}
\flright{\textsc{T.S. Eliot}, \emph{Four Quartets}}
\end{quotex}


Laid up with a little \emph{agita} caused by too much Harp Lager\footnote{\url{https://en.wikipedia.org/wiki/Harp_Beer}} and corned beef\footnote{\url{https://en.wikipedia.org/wiki/New_England_boiled_dinner}}, I got a chance to catch up on some Internet links my ``friends'' have recommended. They turn out to be of a little interest, though probably not in the way they were intended. They are categorized as Excellence, God is Love, and something new under the sun.

The underlying theme is that we already have everything we need. We don't need to explore elsewhere nor seek any modernist or post-modernist novelties.

\paragraph{Seeking Excellence}


A Swami Nirandjananda, speaking in Bulgaria\footnote{\url{https://www.youtube.com/watch?v=McPFZmngym4&feature=share}}, assures us that the true purpose of life is to seek excellence, for which information, he was accorded rousing applause.

Of course, we upholders of the Western tradition recognize this as a call to \emph{arête}\footnote{\url{https://en.wikipedia.org/wiki/Arete}}. Yes, we Westerners are expected to seek excellence, we have known this for quite some time, and we don't, or shouldn't, need an outsider to remind us. Guenon expresses it as the power to actualize the possibilities of the human state, a task naturally suited to the warrior caste. So now Gornahoor urges you all to seek excellence and to actualize your talents. Isn't that enough or are you secretly expecting something more spectacular?

\paragraph{God is Love}


{}\strut \\

An American master of Krishna consciousness also explains the purpose of life in this online lecture\footnote{\url{https://www.bvks.com/2011/07/his-holiness-lectures-at-stanford-university-california}}. He talks about cosmic order, or Logos. He asserts that man is a rational animal. To be human a man must be interested in purposes, or final causes as it is called in the West; otherwise he is no different from an animal. He condemns abortion and homosexuality. Finally, he assures us that God is Love. He claims he was raised as a Roman Catholic, and his talk could have been taken from the Catechism of the Catholic Church.

Of course, he no longer believes in ``that'' God, the sky God who punishes masturbators unjustly. He fails to mention that in the Hindu scheme you may come back as a rutabaga, or worse, as a dalit. What he neglects to mention is that Catholic or Western teaching does not actually teach such a God. The theism of the Western tradition is the God of Being, both transcendent and immanent, and nothing like what was described. Western theism is based on Aquinas, Aristotle, and Plato, with roots deep into paganism. Guenon notes this, and asserts that the Western conception of God is the same as Ishwara in Hinduism\footnote{\url{https://en.wikipedia.org/wiki/Ishwara}}. Now we can fault, perhaps, the priests for not effectively teaching the reality of God, but a public intellectual has an obligation to know better.

This applies \emph{a fortiori} to the theorists of the Nouvelle Droite, who are obsessed with monotheism and its alleged evils. Western theism is nothing like what Mr. Benoist describes, and is certainly not a Semitic notion, as it is based on Roman, Greek, and Indian notions. Apparently, his intuitions of his feelings is not interested in such facts. You can take that critique seriously at your own risk.

\paragraph{Something New under the Sun}


The third link is an interview with the alt right. Whenever I hear this, I am reminded of the Three Stooges'\footnote{\url{https://en.wikipedia.org/wiki/Three_Stooges}} skit on the ``other right'', but I will try to overcome this. Like the New Right, the alt right is based on Nietzsche somehow. About a third of the way into the interview, the question of Nietzsche's influence was posed. To my surprise, the speaker became tongue-tied and was unable to answer the question; I assume because it is impossible to extract a political program from Nietzsche's writings.

Nietzsche taught the will-to-power, not the will-to-knowledge. Hence, there can only be various perspectives, so the search for truth is futile. All the expressions of these perspectives are manifestations of the will to power. It is therefore pointless to dispute their truthfulness. Rather, the approach must be psychological. That is reason that Nietzscheans always attack the motives or psychological health of holders of alternative perspectives. Or else, like New Agers\footnote{\url{https://www.thesecret.tv/}}, they point out the logical consequences of holding a certain perspective. For example, the writers on the new right never ask the question of the truth of the monotheism of the Western tradition, since it is absolutely irrelevant. Rather they point out the alleged consequences of holding that opinion. For example, they claim that a belief in monotheism somehow leads to intolerance or totalitarianism, which are a priori assumed to be (dare I say?) evil.

We hold that the Right can be objectively defined. The Traditional Right is based on the Ancien Regime, of which \textbf{Edmund Burke} and \textbf{Joseph de Maistre} are the exemplars from different perspectives.

Some of them claim that there are some Christian writers, presumably intellectual residues of the Old Right, aligned with the alt right. That seems incredible to me, as the movements are intellectually incompatible. True to his tradition, they are unconcerned with the truth claims of his allies, and attributes psychological motives to them. Namely, they embrace those traditions only because they are dissatisfied with American Protestantism. Gornahoor, on the other hand, recalls Constantine, Charlemagne, Aquinas and Dante; hence we are quite dubious about that assertion.

\url{https://www.youtube.com/watch?v=cx1klDB02bs} 

\flrightit{Posted on 2012-03-19 by Cologero}

\begin{center}* * *\end{center}

\begin{footnotesize}
\begin{sffamily}
\texttt{logres on 2012-03-20 at 00:08 said:}

The title of the interviewing radio show ``Voice of Reason'' reminds me of that propaganda writer in Gironella's trilogy on the Spanish Civil War who was called ``Warning Voice'' -- definitely not portrayed very sympathetically, although plenty of the Requettes and Carlists come out looking just fine.

\url{http://en.wikipedia.org/wiki/Requet\%C3\%A9s}

He looks like an Anglofied Romney.

\hfill

\texttt{Visionsofglory14 on 2012-03-20 at 09:16 said:}

Why does the `Old Right' tend towards greater intolerance for attempts at implementing a New Right than for the formless blob of Mass represented by the modern Left. Why reject an ally to that?

Christian theology teaches: ``God blesses those who love Him, Others have meaningless lives''

Christians want the lives of non-Christian to be meaningless. Where they are not meaningless, they go out of their way to make them meaningless. A non-Christian who has found meaning is a Christian's worst memory.

\hfill

\texttt{Cologero on 2012-03-20 at 13:04 said:}

Here is yet another example of the modern mind at work, much to my disappointment. Once again, rather than addressing the intellectual content, the commentator has to ascribe evil motives to others. Apparently, Visions of glory can read minds and is thus able to ascribe unexpressed desires and intentions to men across the Internet.

An intelligent reader would note these issues:

1) We mentioned Maurras and Evola as examples of the Old Right, neither of whom is a Christian.

2) No one is rejecting allies. We are discussing the roots of ideas. If we have misrepresented Mr. Benoist in any way, then please offer a correction. It seems to use that the New Right is rejecting allies, and for some very unconvincing reasons. \emph{Are we making stuff up? No, we actually quoted Mr. Benoist himself when he specifically and unambiguously rejected Charles Maurras and Julius Evola as guides to the new right}.

3) The post was not written from a Christian of view and nothing was said about the meaningless lives of others. The final two paragraphs of the comment are incomprehensible gibberish.

4) Unfortunately, Visions of glory, despite his education, is apparently unable to grasp a literary allusion. The quote from Elliott was meant to reinforce the idea that meaning can be found at home, after much searching. A careful reading of the post will show that there is no dispute with the swamis per se. Rather, we pointed out that there is no need to turn East since the exact same teaching is found in the West. Was that not totally obvious to anyone else?

\hfill

\texttt{Visionsofglory14 on 2012-03-20 at 16:40 said:}

\url{http://findarticles.com/p/articles/mi_m1141/is_40_36/ai_65643785/}


March 1997: Ratzinger describes Buddhism as ``an auto-erotic spirituality'' in an interview with a French newspaper. Ratzinger said, ``In the 1950s someone said that the undoing of the Catholic church in the 20th century wouldn't come from Marxism but from Buddhism. They were right.''

Quote from China Matters: ``Reportedly, at the time Cardinal Ratzinger was incensed that there were allegedly more Frenchmen studying to be Buddhist monks than Benedictine monks.''

One obvious problem with certain Protestants that makes their view anti-Traditional is their bunker mentality. They see things from a rigidly dualistic , `us versus them' or `Christians versus `the world'' point view that stems straight from the Bible. It cannot form a common front with Tradition because it interprets any creed that is not their own as ``evil,'' ``the world,'' etc. They actually prefer to ally with the forces of global subversion because those forces are easy stereotypes of the the ``evil'' familiar to them in the Bible. Hedonists, gross materialists etc.

Evola referred to those who practice the Catholic Faith in its old form as mere ``half-way Traditionalists.'' (I posit he would dismiss its modern form as anti-Traditional altogether). This means that Catholics who want to form a common front against global subversion must consciously go against the grain with the rest of their global community and even against their Pope. This makes it difficult to find dependable allies within the Church.

It's exceedingly uncommon but not impossible to find good allies within Protestantism now; it's probably about the same for Catholicism.

\hfill

\texttt{Enda Miller on 2012-03-20 at 17:50 said:}

As far as the ``New Right'' is concerned, I have found some of its political insights intriguing and insightful, especially its analysis of classical ``liberalism'' and the dangers of ``communism''. Yet, I have to agree that it's criticism's of Western Monotheism are puerile and unconvincing. For example, Benoist claims, along with many New Right thinkers, that Judeo-Christianity is the father of totalitarianism. However, if one is to be candid, one must admit that totalitarianism is a modern political phenomenon and that the modern era can hardly be defined as Christian. Of course, and perhaps paradoxically, the New Rightists agree with this latter claim, and have shown extensively that all of the modern ``isms'' are the result of the secularization of Christian dogma. Intellectual honesty would suggest that ``New Right'' thinkers are in fact great Inquisitor's, exposing the dangers of leaving Christianity to the Masses, and upholding the neccessity of an all-encompassing Orthodoxy... .;)

\hfill

\texttt{Cologero on 2012-03-20 at 18:00 said:}

Compare this quote from Evola's Men Among the Ruins

\begin{quotex}\footnotesize

Action through ``myths,'' namely through formulas lacking any objective truth and that appeal to the sub-intellectual dimension and passions of individuals and the masses, is the inseparable counterpart of the aforementioned political climate. In the most characteristic modern trends, the notions of ``country'' and ``nation'' display to an eminent degree the quality of myths, susceptible to receiving the most varied contents depending on which way the wind blows and on the political parties, with the only common denominator being the denial of the political principle of pure sovereignty.
\end{quotex}


To what Alain de Benoist has claimed:

1) Unlike Evola, Benoist has clearly articulated his preference for mythos over logos, that is, myth over objective truth

2) Unlike Evola, Benoist has clearly articulated his reliance on the ``intuitions of his passions'', which Evola calls ``sub-intellectual.''

This aligns Benoist with modernity, as someone who denies the principle of pure sovereignty. And Benoist does indeed deny that, and can only see in it a totalitarianism.

So in what conceivable sense can he be considered an ally?

\hfill

\texttt{Cologero on 2012-03-20 at 19:46 said:}

Once again, Visions of Glory, you have posted a random comment with no point and no relation to the post in question.

We are not interested in choosing up sides like some children's game ... check out our post on the Shibboleth to understand why. You have merely given the equivalent of a weather report, describing which way the wind is blowing.

Our main concern is developing the intellectual understanding of Tradition. You seem to be overly concerned with looking over your should to check for ``allies'', usually in the wrong places. That is not the kind of man we are looking for. Create your own allies.

\hfill

\texttt{Cologero on 2012-03-23 at 17:29 said:}

Apparently the Buddhists see things differently: \textit{Just for the Hell of it}\footnote{\url{http://bhikkhublog.blogspot.com/2007/04/just-for-hell-of-it.html}}.

\hfill

\texttt{Cologero on 2012-03-23 at 17:48 said:}

I agree, Mr. Miller, but my intent was not to criticize the New Right per se, but rather to explore the roots of their thought as described by Alain de Benoist himself. Men may agree in certain conclusions, but if their presuppositions differ, then a superficial agreement does not count for much. From the Traditional point of view, what is necessary is the \textit{Common Mind}\footnote{\url{http://www.righthandpath.org/the-church-and-the-common-mind/}}. (Readers may substitute the metaphysical notion of ``Spiritual Authority'' for the theological term ``church''. The point is the same.)

Now, by \emph{polytheism}, the New Rightists do not really mean belief in multiple gods, but rather the assertion that the world consists of multiple powers, independent of each other and irreducible to a common principle. This is metaphysically unacceptable and is really no different from liberalism; instead of individuals themselves being independent monads (as in liberalism), for the new right, different cultures are monads. There is no basis for inter-communication as there is no common mind.

The New Right's claim is based on a historical error. As we pointed out in our several posts on the
\textit{Ancient City}\footnote{\url{http://www.gornahoor.net/?tag=ancient-city}}. The pagans of the Ancient City were not at all tolerant of each other. That is because the cults of different cities were mutually incompatible, so there was all out war between them. Only when they could be arranged hierarchically, could they be at peace with each other. But a hierarchy require a peak ... which leads to monotheism.

\end{sffamily}
\end{footnotesize}

